% ============================================================================
% CHAPTER 2: THEORETICAL BACKGROUND
% ============================================================================

\chapter{Θεωρητικό Υπόβαθρο}
\label{ch:background}

% TODO: Write this chapter based on THESIS_OUTLINE.md Section 2

\section{Αναπαράσταση του Κύβου}
\label{sec:bg_representation}

% Content: Facelet, Cubie, Singmaster notation
% Reference: src/cube/rubik_cube.py, src/kociemba/cubie.py

\subsection{Αναπαράσταση Facelets}

% 54 αυτοκόλλητα, αρίθμηση

\subsection{Αναπαράσταση Cubies}

% 8 γωνίες, 12 ακμές, προσανατολισμός

\subsection{Συμβολισμός Singmaster}

% U, D, L, R, F, B και παραλλαγές


\section{Στοιχεία Θεωρίας Ομάδων}
\label{sec:bg_group_theory}

% Content: Groups, permutations, subgroups, cosets
% Reference: docs/notes/01_group_theory_fundamentals.md


\section{Αλγόριθμοι Αναζήτησης}
\label{sec:bg_search}

\subsection{Αλγόριθμος A*}

% f(n) = g(n) + h(n), admissibility

\subsection{Αλγόριθμος IDA*}

% Iterative deepening, memory efficiency
% Reference: src/korf/a_star.py

\subsection{Pattern Databases}

% Concept, admissibility, additive heuristics
% Reference: src/korf/pattern_database.py
